\section*{Chapter \ref{ch:rs:fundamental-analyst-and-event-features} --- Fundamental, Analyst and Event Features}

\begin{solution}{exo:rs:fundamental-analyst-and-event-features:1}
$\$4.80 - \$3.30 = \$1.50$. It is first known when the 10-K is filed, 55 days after 30 June: 24 August (or at the earnings release, if
the store keys on releases).
\end{solution}

\begin{solution}{exo:rs:fundamental-analyst-and-event-features:2}
Forty days after 30 September: 9 November. Sixty days after 31 December in a year of 365 days: 1 March. (A deadline that falls on a weekend
or holiday moves to the next business day.)
\end{solution}

\begin{solution}{exo:rs:fundamental-analyst-and-event-features:3}
$1.42/0.355 = 4$: a four-for-one split rescaled the history. The store keeps \$1.42 with its first filing date as \omterm{def:rs:point-in-time-data-and-the-biases:bitemporal}{knowledge time} and \$0.355
with the later filing's, and records the split as an event with its own date; a query rescales to the share basis of its decision date.
\end{solution}

\begin{solution}{exo:rs:fundamental-analyst-and-event-features:4}
The seasonal differences are 0.04, 0.02, 0.06, 0.02, 0.02, 0.04, $-0.02$ and 0.18. The six before the last (0.02, 0.06, 0.02, 0.02, 0.04,
$-0.02$) have a sample standard deviation of 0.0266, so the last quarter's SUE is $0.18/0.0266 = 6.8$.
\end{solution}

\begin{solution}{exo:rs:fundamental-analyst-and-event-features:5}
The window from day 1 to day 21 after the quarter's end contains the announcement when it falls on one of those 21 of the 63 possible days: a
third (0.330 measured). In those windows the feature sees the announcement's jump, whose size dwarfs the drift the correct key measures; a third
of the windows carrying a jump several times larger than a quarter's drift is enough to multiply the IC by eight.
\end{solution}

\begin{solution}{exo:rs:fundamental-analyst-and-event-features:6}
A look-ahead gains what the price does between the period end and the \omterm{def:rs:point-in-time-data-and-the-biases:bitemporal}{knowledge time} because of the hidden number. An \omterm{def:rs:fundamental-analyst-and-event-features:sue}{earnings surprise} moves the
price on its release, inside the gap; a book value's filing moves it little, and book-to-price changes slowly anyway, so the window's return is
nearly independent of which of the two keys is used. Book-to-price would inflate if the book value carried news the price reacted to on the filing
(a large write-down), or if the period-end price used in the ratio were replaced by a later one (a price look-ahead).
\end{solution}

\begin{solution}{exo:rs:fundamental-analyst-and-event-features:7}
\texttt{rs\_fundamentals.release\_to\_filing}: Microsoft filed every 10-Q on the release day, Procter \& Gamble 95\% of them; Coca-Cola 7\% and
Intel 10\%. For the others, the median days from release to filing are one (Apple, Intel), two (Coca-Cola, Procter \& Gamble), five (Exxon Mobil),
seven (Home Depot), 14 (Johnson \& Johnson, Nike) and 16 (Walmart).
\end{solution}

\begin{solution}{exo:rs:fundamental-analyst-and-event-features:8}
It is the most accurate history and the wrong one for a backtest: every value is today's, so a decision in the past sees \omterm{def:rs:point-in-time-data-and-the-biases:vintage}{restatements} (Microsoft's
fiscal 2017 EPS of \$3.25 instead of the \$2.71 known until August 2018) and split rescalings it could not have known. Research needs the values as
first reported, with every later vintage and its date.
\end{solution}

\begin{solution}{pb:rs:fundamental-analyst-and-event-features:1}
\begin{enumerate}
\item 542 period ends; 541 are followed by a release within 90 days.
\item Quarters: 23 days to the release, 29 to the 10-Q. Fiscal years: 26 to the release, 48 to the 10-K.
\item 41 and 60 days; the deadlines are 40 (moved to the next business day when it falls on a holiday or weekend) and 60 days for large filers.
\item 16\% on the same day; the median gap is 7 days.
\item 31, of which 24 by an exact ratio.
\item Seven-for-one on 6 June 2014 and four-for-one on 28 August 2020; per-share values before each were later reported divided by 7 and by 4 in
the filings that repeated them.
\item First reported at \$2.71, restated to \$3.25 in the fiscal 2018 filings: Microsoft adopted a new revenue standard with the full
retrospective method, which restated each prior period presented.
\item \$28.82 from first-reported quarters (two before the split, two after), against an annual EPS of \$6.45.
\item As first reported with their \omterm{def:rs:point-in-time-data-and-the-biases:bitemporal}{knowledge times}, with splits stored as events, and rescaled at query time to the decision date's share
basis.
\item 31 trading days, about six weeks.
\item 0.020 (standard error 0.006) keyed on the announcement; 0.167 (0.006) keyed on the period end.
\item A third (0.330).
\item 0.029 (0.038) keyed on the period end; $-0.005$ (0.023) keyed on the filing, with the period-end price.
\item Each quarter's IC is dominated by the value factor's return in its window, and nothing happens on the filing day to make the two windows'
returns depend on the ratio differently: the expected inflation is small and the noise large.
\item \textbf{Named result.} In \texttt{firm.synthmkt}, keying the \omterm{def:rs:fundamental-analyst-and-event-features:sue}{earnings surprise} on the period end instead of its announcement multiplies
its IC by eight (0.020 to 0.167); keying book-to-price on the period end instead of its filing changes it by less than its standard error.
\item The days between the release and the filing (a week at the median in the EDGAR data): for a drift feature, the start of the drift.
\item The period-end key, by far, for a SUE feature: it trades the announcement itself. A restated history distorts the size of some surprises
but not their timing.
\item Recompute the feature's IC with the entry shifted to the \omterm{def:rs:point-in-time-data-and-the-biases:bitemporal}{knowledge time}; plot the IC against the shift, since a leak shows as a
cliff at the true publication lag; check the timing of the returns earned around announcements.
\item Every value with its period, its first publication time, every later vintage with its time, and corporate actions as dated events.
\item The time it was first published, at the release or, for a database fed by filings, at the filing, never the end of the period it describes.
\end{enumerate}
\end{solution}

\begin{solution}{iq:rs:fundamental-analyst-and-event-features:1}
Data stored with the time each value became known, and queried as of a decision time, so a backtest sees only what was known then. It matters
because accounting numbers are published weeks after their periods and are later restated or rescaled; a database of current values keyed on
period ends leaks both.
\end{solution}

\begin{solution}{iq:rs:fundamental-analyst-and-event-features:2}
Difference the year-to-date figures: the second quarter is the six months minus the first, the fourth the year minus the nine months; each derived
quarter is known when its later component is. Sum four consecutive quarters as known at the decision date, on one share basis for per-share figures.
\end{solution}

\begin{solution}{iq:rs:fundamental-analyst-and-event-features:3}
The tendency of returns after an earnings release to continue in the direction of the surprise (Foster, Olsen and Shevlin, 1984). Feature: SUE
against a seasonal random walk or the consensus, keyed on the release, ranked among the stocks that announced recently, held for weeks; check
the key date and neutralise size.
\end{solution}

\begin{solution}{iq:rs:fundamental-analyst-and-event-features:4}
The filing date is right: the period end uses numbers before they were known. The improvement is large when the number's publication moves the
price (earnings) and small when it does not (a slowly changing book value); on the chapter's simulated market book-to-price moved by less than its
standard error, the surprise by a factor of eight. A third is a warning either way.
\end{solution}

\begin{solution}{iq:rs:fundamental-analyst-and-event-features:5}
Timestamps (forecast date against the date the vendor recorded it), revisions that overwrite history, survivorship of analysts and of covered firms,
coverage chosen on expected performance, stale forecasts in the consensus, and optimism that varies with the analyst's incentives.
\end{solution}

\begin{solution}{iq:rs:fundamental-analyst-and-event-features:6}
A \omterm{def:rs:point-in-time-data-and-the-biases:bitemporal}{bitemporal} store of facts keyed by (entity, field, period) with every vintage and its \omterm{def:rs:point-in-time-data-and-the-biases:bitemporal}{knowledge time}, values kept as reported on their own share
basis; a separate table of split events with their dates; a query (entity, field, period, decision time) that takes the last vintage known then and
multiplies it by the product of the split factors between the vintage's basis and the decision date's, both taken as known at the decision time.
\end{solution}
